\section{Discussion and conclusion}\label{sec:discussion}
The fitted parameters $\theta$ observe the structure in modal space: poles,
rank-one residues and the port rows of the modal matrix, and with them the
modal displacement under the measured forces. The nominal model then embeds
this fixed $r$-mode model in the mesh. The embedding keeps the measured
ports exactly and every modal coordinate with a margin, and within these
constraints it minimizes the nominal $F=ma$ residual along the measured,
damped motion. The damping is never compared with the nominal model, which
has none; it enters through the motion. The realized operators satisfy the
fitted modal equations exactly, even where the nominal operators do not.

The chain shows the mechanism. The fit fixes two endpoint rows, the
completion fills the middle row, and the equality-only map already meets
the margin. The laboratory shell shows the limits. The fitted transfers are
preserved to roundoff, and the constrained iteration converges with a
margin cost of \LabMarginCost\% in $J_F$. But $J_F$ constrains only the
motion that the test produces. With two shakers, correlated modes leave
combinations that the criterion barely sees. The completion gives them
large amplitudes, and the excluded responses are predicted worse than with
the undamped criterion: \LabExcludedError\% against
\LabExcludedErrorJZero\%. These assessments selected neither the completion
nor its margin. More independent inputs, or a weight that also covers the
unexcited combinations, would constrain these directions. Forces at known
positions are recovered with errors of the order of the port misfit. The
saved factors, operators, spectra and source reproduce every calculation
and figure.
